\providecommand{\main}{..}
\documentclass[\main/main.tex]{subfiles}
\begin{document}

La teoria MAUT (multiple attribute utility theory) ipotizza di trovarsi nella situazione in cui un decisore è in grado di ordinare gli indicatori, creando un set $\Pi$ con un ordine debole. Ipotizza inoltre l'esistenza di una funzione valore $u(f)$ conforme a $\Pi$. Spesso la $u(f)$ non si conosce ed è necessario costruire il passaggio da $\pi$ a $u(f)$

\begin{figure}[H]
  \centering
  \includegraphics[width=0.75\textwidth]{\main/images/MAUT.png}
  \caption{MAUT in progress}
\end{figure}

Campiono quuindi gli impatti e connetto quelli che risultano indifferenti preferenzialmente, cercando di identificare forme analitiche che le descrivano.

\[
  u(b) = f_1f_2 \rightarrow u(b) = f_1^{\alpha_{1}}f_2^{\alpha_{2}}
\]

Definita quindi una formula vado a verificare che valga su quei punti:

\[
  u(f_1) = u(f_2) \Rightarrow  f_1^{\alpha_{1}}f_1^{\alpha_{2}} = f_2^{\alpha_{1}}f_2^{\alpha_{2}}
\]

\begin{figure}[H]
  \[
    u(f) = \prod_{l=1}^p f_l^{\alpha_l}
  \]
  \caption{\textbf{Funzioni di Cobb-Douglas:} descrive un consumatore, viene usata in economia}
\end{figure}

Questo procedimento diventa rapidamente insolubile all'aumentare dei decisori e degli indicatori. Viene quindi \textit{ipotizzato} che la funzione di utilità sia \textbf{additiva}, cioè che valga:

\[
  u(f) = \sum_{l=1}^p u_l(f_l)
\]

\section{Indipendenza preferenziale}
Dato un insieme $P = \{1, \dots, p\}$ degli indici, $L \subset P$ con $L$ preferenzialmente indipendente da $P \setminus L$.

\subsubsection{Esempio sui viaggi}
Viaggiare bene prima implica voler viaggiare bene dopo?
\[
  \begin{bmatrix}
    f_L \\
    f
  \end{bmatrix}
  \preceq
  \begin{bmatrix}
    f'_L \\
    f
  \end{bmatrix}
  \Leftrightarrow
  \begin{bmatrix}
    f_L \\
    \xi
  \end{bmatrix}
  \preceq
  \begin{bmatrix}
    f'_L \\
    \xi
  \end{bmatrix}
  \quad
  \forall f_L \forall f'_L \forall f \forall \xi
\]

\subsubsection{Esempio sui pranzi}
Primo e secondo non sono preferenzialmente indipendenti.

\[
  f = \begin{bmatrix}
    primo \\
    secondo
  \end{bmatrix}
\]

\end{document}